% GENEROVÁNO skriptem doc/tools/md2tex.py ze souboru kapitoly/04_navrh_a_metodika.md. Needitovat, upravuje se Markdown.

% Kapitola 4, verze 2 (2. 11. 2026): bez poznatků z testování (ty jsou v poznamky_k_prenosu.md), přepsáno přímějším stylem. Převod do LaTeXu: doc/tools/md2tex.py.

\chapter{Návrh a~metodika srovnání}
\label{kap:metodika}

Ke každému algoritmu z~této práce vznikly tři implementace: v~jazyce C, ve Vitis \ac{HLS} a~ve \ac{VHDL}. Aby jejich porovnání mělo smysl, musí platit dvě věci. Všechny tři verze musí počítat přesně totéž a~všechny se musí měřit stejným způsobem. Tato kapitola popisuje, jak jsem obojí zajistil. Naměřená čísla jsou až v~kapitole~\ref{kap:vysledky}.

\section{Cíle a~kritéria srovnání}
\label{sec:kriteria}

Zadání po mně chce implementovat vybrané algoritmy v~softwaru a~hardwaru, a~to ve VHDL i~pomocí vysokoúrovňové syntézy, a~hardwarové verze porovnat se softwarem. Každý algoritmus proto implementuji třikrát se stejnou aritmetikou: v~jazyce C, který je zároveň referenčním modelem a~softwarovým řešením pro procesor Cortex-A53, ve Vitis HLS a~ručně ve VHDL. Obě hardwarové verze mají stejné rozhraní a~stejnou mapu registrů, takže jsou zaměnitelné a~ovládá je jeden ovladač v~jazyce C.

Na prvním místě je správnost. Implementace, která se s~referenčním modelem neshoduje bit po bitu, se dál neposuzuje, protože porovnávat rychlost nebo velikost špatně počítajícího obvodu nedává smysl. Správné implementace porovnávám podle kritérií v~tabulce~\ref{tab:kriteria}. Prostředky jsou počty buněk \ac{LUT} a~\ac{FF}, bloků \ac{DSP} a~pamětí \ac{BRAM}, které jádro v~obvodu zabere. Prostředky a~maximální frekvenci měřím na samotném jádře, propustnost a~rychlost proti softwaru na celém řetězci na desce.

\begin{table}[htbp]
    \centering
    \caption{Kritéria srovnání HLS, VHDL a~softwaru}
    \label{tab:kriteria}
    \begin{tabular}{p{3.3cm}p{5.2cm}p{5.2cm}}
        \toprule
        Kritérium & Co se měří & Jak a~kde \\
        \midrule
        Správnost & počet chybných bajtů proti referenčnímu modelu & simulace a~měřicí aplikace na desce (oddíl~\ref{sec:overeni}) \\
        Prostředky & \acs{LUT}, \acs{FF}, \acs{DSP} a~\acs{BRAM} jednoho jádra & syntéza mimo kontext (oddíl~\ref{sec:ooc}) \\
        Maximální frekvence & nejkratší perioda hodin, kterou jádro zvládne & syntéza mimo kontext (oddíl~\ref{sec:ooc}) \\
        Propustnost & doba zpracování snímku 1920\,×\,1080 bodů v~ms a~počet snímků za sekundu & měření na desce (oddíl~\ref{sec:deska}) \\
        Rychlost proti softwaru & poměr doby softwarového a~hardwarového zpracování stejného algoritmu & měření na desce, software na procesoru Cortex-A53 (oddíl~\ref{sec:sw}) \\
        \bottomrule
    \end{tabular}
\end{table}

\section{Referenční model v~jazyce C}
\label{sec:model}

Jako referenci bych mohl použít přímo OpenCV, jenže ten má pro můj účel jednu nevýhodu: část operací počítá v~plovoucí řádové čárce nebo zaokrouhluje jinak, než by to udělal obvod. Hardware potřebuje definici, kterou lze bit po bitu zopakovat. Napsal jsem proto vlastní celočíselnou knihovnu v~jazyce C (\texttt{lib\_imgproc}). Kde to jde, drží se chování OpenCV. Zaokrouhlení, saturaci a~ošetření okrajů ale stanovuje výslovně.

Model má dvě role. Je měřítkem správnosti hardwarových verzí a~zároveň je softwarovým řešením, které se na procesoru Cortex-A53 měří proti hardwaru.

Konstanty, ve kterých by se verze mohly rozejít, neopisuji ručně. Generátor z~definic vytvoří koeficienty převodu barev a~tabulky pro převod mezi RGB a~HSV a~zapíše je do hlavičkového souboru pro C a~HLS a~do balíčku pro VHDL. Všechny tři implementace tak používají stejné hodnoty.

Protože je model měřítkem pro všechno ostatní, porovnávám nejdřív jeho samotný s~OpenCV. Odchylky od OpenCV jen vyhodnotím (kapitola~\ref{kap:vysledky}). Hardwarové verze se vždy porovnávají s~modelem, ne s~OpenCV. Aritmetiku, převody barev, medián a~filtry kontrolují navíc jednotkové testy, které běží i~na procesoru ARM.

% TODO(citace): dokumentace OpenCV (cvtColor, medianBlur, threshold, filter2D) jako zdroj sémantiky operací; podklad k vložení do citace.bib dodá Karel z Citace PRO.

\section{Aritmetika s~pevnou řádovou čárkou}
\label{sec:aritmetika}

Obrazové body jsou osmibitová celá čísla, ale koeficienty převodu barev nebo měřítko filtru jsou zlomky. Počítat je v~plovoucí řádové čárce by v~obvodu \ac{FPGA} stálo víc prostředků a~trvalo by to déle než celočíselné počítání. Bloky \ac{DSP} mají navíc celočíselnou násobičku se vstupy širokými 27 a~18 bitů. Zlomky proto počítám v~pevné řádové čárce.

Číslo ve formátu $Q_{m.n}$ je celé číslo se znaménkovým bitem, $m$ bity celé části a~$n$ zlomkovými bity. Uložené celé číslo $X$ představuje hodnotu

\begin{equation}
    x = \frac{X}{2^{n}} .
    \label{eq:q-format}
\end{equation}

Pro převody barev jsem zvolil formát $Q_{2.15}$. Má 18 bitů, rozsah přibližně od $-4$ do $4$, což pro koeficienty stačí, a~rozlišení $2^{-15}$. Osmnáct bitů přesně vyplní širší vstup násobičky v~bloku DSP48E2. Koeficient 0,299 se v~tomto formátu uloží jako celé číslo $9\,798$, protože $0{,}299 \cdot 2^{15} = 9\,797{,}6$.

Pevná řádová čárka vyžaduje dvě konvence, které jsou ve všech třech implementacích stejné. První je zaokrouhlení polovin nahoru při posunu o~$n$ bitů vpravo, druhá saturace výsledku na rozsah osmibitového obrazového bodu:

\begin{equation}
    \mathrm{round\_shift}(X, n) = \left\lfloor \frac{X + 2^{n-1}}{2^{n}} \right\rfloor , \qquad \mathrm{sat}_{u8}(v) = \min\bigl(\max(v, 0), 255\bigr) .
    \label{eq:round-sat}
\end{equation}

V~hardwaru je dělení mocninou dvou jen aritmetický posun, takže zaokrouhlení stojí jednu sčítačku, a~u~bloku DSP48E2 dokonce žádnou, protože se provede v~jeho závěrečné sčítačce.

Celý postup je vidět na příkladu. Řádek matice převodu RGB na šedou (BT.601) má v~pořadí složek, v~jakém leží v~paměti, koeficienty $19\,235$ (zelená), $3\,735$ (modrá) a~$9\,798$ (červená). Vstupem je pixel s~červenou 200, zelenou 100 a~modrou 50. Součet součinů je $19\,235 \cdot 100 + 3\,735 \cdot 50 + 9\,798 \cdot 200 = 4\,069\,850$. K~němu se přičte polovina, tedy $2^{14} = 16\,384$, což dává $4\,086\,234$. Po posunu o~15 bitů je výsledkem $124$, protože hodnota před posunem je 124,7 a~posun zaokrouhlí dolů. Výpočet v~plovoucí řádové čárce dává $0{,}299 \cdot 200 + 0{,}587 \cdot 100 + 0{,}114 \cdot 50 = 124{,}2$, tedy také 124.

Koeficienty jednoho řádku matice musí dát v~součtu přesně $2^{15}$, aby se bílá barva převedla na bílou. Prosté zaokrouhlení každého koeficientu zvlášť to nezaručuje. Generátor proto koeficienty zaokrouhlí dolů a~chybějící jednotky přidělí těm koeficientům, které měly největší zlomkovou část (metoda největšího zbytku).

Šířky mezivýsledků nikde nepíšu jako konstanty. Součin dvou čísel šířky $W_a$ a~$W_b$ má $W_a + W_b$ bitů a~součet $N$ čísel šířky $W$ má $W + \lceil \log_{2} N \rceil$ bitů. Jádra jsou díky tomu generická: například konvoluce se pro jiný rozměr jádra nepřepočítává ručně, a~nehrozí tak chyba v~šířce při změně parametru. Formáty a~konvence jednotlivých operací shrnuje tabulka~\ref{tab:formaty}.

\begin{table}[htbp]
    \centering
    \caption{Formáty a~konvence aritmetiky u~jednotlivých operací}
    \label{tab:formaty}
    \begin{tabular}{p{3.3cm}p{5.2cm}p{5.2cm}}
        \toprule
        Operace & Formát a~měřítko & Zaokrouhlení a~saturace \\
        \midrule
        Převod barev (matice 3\,×\,3) & koeficienty $Q_{2.15}$, offset se přičítá v~akumulátoru před posunem & zaokrouhlení polovin nahoru, saturace 0 až 255 \\
        RGB na HSV & celočíselný algoritmus OpenCV, posun o~12 bitů, tabulky podílů a~odstínu, odstín v~rozsahu 0 až 179 & shodné s~OpenCV \\
        HSV na RGB & vlastní celočíselný algoritmus s~tabulkami převrácených hodnot & zaokrouhlení polovin nahoru, saturace \\
        Konvoluce $K \times  K$ & osmibitové celočíselné jádro, přesná akumulace, měřítko $Q_{.14}$ a~posun hodnoty \texttt{delta} & zaokrouhlení polovin nahoru, saturace \\
        Medián, prahování, histogram & jen celá čísla (porovnávání a~počítání) & není potřeba \\
        \bottomrule
    \end{tabular}
\end{table}

% TODO(citace): Q formát / aritmetika s pevnou řádovou čárkou (např. Texas Instruments, ARM, MathWorks), z původních poznámek v DP.tex; podklad k vložení dodá Karel.

\section{Úrovně ověření}
\label{sec:overeni}

Bitovou shodu ověřuji na třech úrovních, které se doplňují. První, referenční model proti OpenCV, je popsaná v~oddílu~\ref{sec:model}. Druhá a~třetí ověřují hardwarová jádra proti modelu (tabulka~\ref{tab:urovne}).

\begin{table}[htbp]
    \centering
    \caption{Úrovně ověření hardwarových jader}
    \label{tab:urovne}
    \begin{tabular}{p{2.4cm}p{3.6cm}p{3.4cm}p{3.7cm}}
        \toprule
        Úroveň & Co ověřuje & Nástroj & Kritérium \\
        \midrule
        HLS & kód C++ a~RTL vygenerované z~něj & C-simulace a~kosimulace C/RTL ve Vitis HLS & shoda dat i~signálů \texttt{TUSER} a~\texttt{TLAST} \\
        VHDL & ručně psaný RTL kód & simulace v~xsim (Vivado) a~GHDL s~generickým testbenchem & shoda dat i~signálů \texttt{TUSER} a~\texttt{TLAST} \\
        Deska & celý řetězec s~procesorem, pamětí a~DMA & měřicí aplikace na procesoru Cortex-A53 & nulový počet chybných bajtů proti modelu a~nulový stavový registr \\
        \bottomrule
    \end{tabular}
\end{table}

Vstupní a~očekávané vektory vytváří nástroj z~referenčního modelu, takže se nic neopisuje ručně. U~HLS ověřuje C-simulace kód v~C++ a~kosimulace stejný test nad RTL, které z~něj Vitis HLS vytvořil. Testbench jádra posílá data v~datovém toku \acs{AXI}4-Stream, ve kterém \texttt{TUSER} označuje začátek snímku a~\texttt{TLAST} konec řádku. Zdroj vkládá do toku náhodné mezery a~příjemce náhodně brzdí, aby se ověřilo, že se žádný pixel neztratí ani nezdvojí a~že platí zpětné řízení toku. Každý test běží nejméně na dvou snímcích za sebou, protože jádro nesmí držet stav z~předchozího snímku. Převod RGB na HSV a~zpět navíc ověřuji v~simulaci na všech $2^{24}$ vstupech.

\section{Metodika měření}
\label{sec:metodika-mereni}

\subsection{Syntéza mimo kontext}
\label{sec:ooc}

Prostředky a~maximální frekvenci měřím syntézou mimo kontext (out-of-context, \acs{OOC}). Jádro se syntetizuje, rozmístí a~propojí samo, bez zbytku návrhu (DMA, propojení AXI, procesor). Porty se nepřipojují na piny a~jediným omezením jsou hodiny s~periodou 5\,ns, tedy 200\,MHz. Výsledek tak nezávisí na tom, co leží vedle jádra: v~celém návrhu by rozmístění a~časování jádra ovlivňovalo, jak je zbytek obvodu zaplněný.

Aby bylo srovnání HLS a~VHDL férové, procházejí obě verze stejným postupem. RTL vygenerované z~HLS jde do stejného skriptu jako VHDL, pro stejný obvod (xck26-sfvc784-2LV-c) a~se stejným cílem 5\,ns. Používám Vivado 2025.2 a~jen čísla po implementaci, ne odhady z~HLS syntézy.

Maximální frekvenci vyjadřuji nejkratší periodou hodin, kterou jádro zvládne:

\begin{equation}
    T = \max\left(5 - \mathrm{WNS},\; 5 - \mathrm{WPWS}\right) \quad [\mathrm{ns}] .
    \label{eq:perioda}
\end{equation}

\ac{WNS} je časová rezerva nejhorší logické cesty a~\ac{WPWS} je rezerva minimální šířky pulzu hodin. Druhý člen je nutný, protože bloky BRAM a~DSP mají minimální periodu danou výrobcem (3,077\,ns u~BRAM a~3,262\,ns u~DSP) a~žádná úprava logiky ji nezkrátí. Vlastní export nástroje Vitis HLS počítá periodu jen z~WNS, a~proto bych jím u~jader s~BRAM nebo DSP dostal kratší periodu, než obvod skutečně zvládne. Počet bloků BRAM udávám v~dlaždicích o~velikosti 36\,Kb.

\subsection{Měření na desce}
\label{sec:deska}

Snímek 1920\,×\,1080 bodů s~24 bity na pixel zabírá 6,2\,MB. Při jednom pixelu za takt a~hodinách 200\,MHz projde obvodem za $1920 \cdot 1080 / 200\,\mathrm{MHz} = 10{,}368\,\mathrm{ms}$, tedy přibližně 96 snímků za sekundu. To je strop, který žádné jádro nepřekoná, a~proto ho měřím jako první, na průchodu bez jádra (bypass). Rozdíl proti průchodu pak říká, kolik času jádro přidává.

Propustnost měřím na desce KV260 v~testovacím řetězci, ve kterém je všech 14 jader v~jednom bitstreamu na 200\,MHz (podrobně oddíl~\ref{sec:retezec}). Procesor nastaví registry vybraného jádra a~\ac{VDMA} přečte snímek z~paměti, pošle ho jádrem a~výsledek zapíše zpět do paměti. Měří se tedy celý řetězec paměť, jádro, paměť, ne samotné jádro.

Testovací snímek aplikace při startu vytvoří deterministicky tak, aby obsahoval plynulé přechody (pro převody barev a~HSV), ostré hrany (pro Sobelův filtr) a~šum (pro medián). Doba zpracování je čas od spuštění čtení z~paměti do jeho zastavení po 20 snímcích. Během přenosu se čte jen stav tohoto kanálu, jednou za 100\,µs, což odpovídá nejistotě 5\,µs na snímek, tedy méně než 0,1\,\%. Měření se opakuje pětkrát a~výsledkem je medián.

Po měření aplikace porovná výstup bajt po bajtu s~referenčním modelem přeloženým pro ARM. Měření je úspěšné, jen když je počet chybných bajtů nulový a~u~jader s~okénkem je nulový i~stavový registr, který hlásí poruchu synchronizace řádků. Výsledky aplikace vypisuje na sériovou linku a~automatický skript je ukládá ve formátu CSV.

\subsection{Softwarové řešení na procesoru Cortex-A53}
\label{sec:sw}

Referenční model přeložený pro ARM spouštím na jednom jádru procesoru Cortex-A53 s~hodinami 1333\,MHz (podle nastavení generovaného ve Vivadu), bez operačního systému a~se zapnutými vyrovnávacími pamětmi. Neběží tak žádný plánovač, který by měření rušil. Kód překládám s~optimalizací \texttt{-O0}, \texttt{-O2} i~\texttt{-O3} překladačem dodávaným s~nástrojem Vitis 2025.2, protože na úrovni optimalizace závisí, jak výhodné porovnání s~hardwarem vychází. Čas se měří čítačem procesoru a~výsledkem je medián z~10 běhů každého algoritmu na stejném testovacím snímku jako u~hardwaru.
